\begin{man}[stdio]{printf}
    \begin{synopsis}
        \begin{plaincode}[gobble=8]
        #include <stdio.h>
        int printf(const char *format, ...);
        \end{plaincode}
    \end{synopsis}

    \begin{desc}
        \codeinline{printf()} schreibt formattierten Text in die Standardausgabe (stdout). Der \codeinline{format}-String enthält:
        \begin{itemize}
            \item Normale Zeichen (werden direkt ausgegeben)
            \item Format-Specifiers (beginnen mit \codeinline{\%})
        \end{itemize}

        \textbf{Die wichtigsten Format-Specifiers:}

        \begin{tabular}{ll}
        \texttt{\%d} & Ganzzahl (dezimal) \\
        \texttt{\%f} & Gleitkommazahl \\
        \texttt{\%c} & Zeichen (char) \\
        \texttt{\%s} & Zeichenkette (String) \\
        \texttt{\%x} & Hexadezimal (klein) \\
        \texttt{\%X} & Hexadezimal (groß) \\
        \texttt{\%p} & Zeiger (Pointer) \\
        \end{tabular}
    \end{desc}

    \begin{example}
        \begin{code}
#include <stdio.h>

int main() {
    int zahl = 42;
    printf("Die Zahl ist: %d\n", zahl);
    printf("Adresse: %p\n", (void*)&zahl);
    return 0;
}
        \end{code}
    \end{example}

    \begin{returns}
        Die Anzahl der ausgegebenen Zeichen, oder \codeinline{-1} bei Fehler.
    \end{returns}

    \begin{seeother}
        \codeinline{scanf()}, \codeinline{sprintf()}, \codeinline{fprintf()}
    \end{seeother}
\end{man}

